\documentclass{article}
\usepackage[utf8]{inputenc}
\usepackage[T1]{fontenc}
\usepackage{amsmath, amssymb, amsthm}
\usepackage{geometry}
\usepackage{hyperref}
\usepackage{bookmark}
\usepackage{tikz}

\usepackage{titling}

% \pretitle{\begin{center}\Huge\bfseries}
% \posttitle{\end{center}}
% \preauthor{\begin{center}\large}
% \postauthor{\end{center}}
% \predate{\begin{center}\large}
% \postdate{\end{center}}

% \renewcommand{\thesubsection}{\Roman{subsection}}
% \newcommand{\exo}[2]{\textbf{\Roman{#1} #2} }

% \geometry{left=2.5cm, right=2.5cm, top=2.5cm, bottom=2.5cm}

\newcommand{\subtitle}[1]{
    \posttitle{\par\end{center}\begin{center}\large#1\end{center}\vskip0.5em}
}


\setlength{\parindent}{0pt}
\setlength{\parskip}{0.8em}

\title{Equation différentielle ordinaire}

\subtitle{Methodes de résolution numérique} 

\author{Axel Lavigne}
\date{\today}


\newcommand{\exo}[1]{
    \vspace{5mm}
    \textbf{Exercice #1}
    \hrulefill
}

\newcommand{\smalltext}[1]{
    \scriptsize
    \textit{#1}
    \normalsize
}

\begin{document}

\maketitle

\section*{Equation différentielle linéaires}

\begin{equation*}
    y^{(n)} + a_{n-1}y^{(n-1)} + \ldots + a_1y' + a_0y = b(x)
\end{equation*}

\textit{methode de résolution générale:}
\begin{itemize}
    \item[--] Solution de l'équation homogène: $y_h(x)$
    \item[--] Solution particulière: $y_p(x)$
    \item[--] Solution générale: $y(x) = y_h(x) + y_p(x)$
\end{itemize}

\section*{Résoudre l'équation homogene}

\hrulefill \,
\textbf{Ordre 1}
\hrulefill 

\begin{equation*}
    y_h' + a(t)y_h = 0
\end{equation*}


\begin{center}
    \textit{Solution } \\ \vspace{1mm}
    \boxed{
        y_h(t) = \lambda e^{-\!\!\int\!a(t)\,dt}
    }
\end{center}

\hrulefill \,
\textbf{Ordre 2}
\hrulefill 
\begin{center}
    Dans le cas ou les coefficients $a$ et $b$ sont constants
    \begin{equation*}
        y_h'' + ay_h' + by_h = 0
    \end{equation*}

    Equation caractéristique
    \begin{equation*}
        r^2 + ar + b = 0
    \end{equation*}
    \begin{equation*}
        \Delta = a^2 - 4b
    \end{equation*}

    % On distingue 3 cas:
    % \begin{itemize}
    %     \item[--] $\Delta > 0$ : 2 racines réelles
    %     \item[--] $\Delta = 0$ : 2 racines réelles égales
    %     \item[--] $\Delta < 0$ : 2 racines complexes conjuguées
    % \end{itemize}
    \hrulefill \,
    \textbf{$\Delta > 0$}
    \hrulefill
    \begin{equation*}
        r_1 = \frac{-a + \sqrt{\Delta}}{2} \quad r_2 = \frac{-a - \sqrt{\Delta}}{2}
    \end{equation*}
    \textit{Solution } \\ \vspace{1mm}
    \boxed{
        y_h(t) = \lambda e^{r_1 t} + \mu e^{r_2 t}
    }

    \vspace{4mm}

    \hrulefill \,
    \textbf{$\Delta = 0$}
    \hrulefill
    \begin{equation*}
        r = -\frac{a}{2}
    \end{equation*}
    \textit{Solution } \\ \vspace{1mm}
    \boxed{
        y_h(t) = (\lambda x + \mu)e^{rt}
    }

    \vspace{4mm}

    \hrulefill \,
    \textbf{$\Delta < 0$}
    \hrulefill
    \begin{equation*}
        r = -\frac{a + i\sqrt{|\Delta|}}{2} \quad 
    \end{equation*}
    \begin{equation*}
        \alpha = -\frac{a}{2} \quad \beta = \frac{\sqrt{|\Delta|}}{2}
    \end{equation*}
    \textit{Solution } \\ \vspace{1mm}
    \boxed{
        y_h(t) = e^{\alpha t}(\lambda \cos(\beta t) + \mu \sin(\beta t))
    }

    % \textit{Solution } \\ \vspace{1mm}
    % \boxed{
    %     y_h(t) = e^{-\frac{a}{2}t}(\lambda \cos(\frac{\sqrt{|\Delta|}}{2}t) + \mu \sin(\frac{\sqrt{|\Delta|}}{2}t))
    % }
    
\end{center}

\section*{Trouver la solution particulière}

La methode générale est de trouver un fonction de la forme de $b(x)$
par exemples si $b(x)$ est un polynôme, on cherche une solution de la forme d'un polynôme

\textit{Note:} Si $b(x)$ est une combinaison de fonctions, on utilise la méthode de superposition:
on trouve une solution particulière pour chaque fonction et on somme les solutions

\hrulefill \,
\textbf{Ordre 1}
\hrulefill

plusieurs cas en fonction de $b(x)$:
\begin{itemize}
    \item[--] $b(x) $ constant $ \longrightarrow $ solution constante
    \item[--] $b(x) $ polynôme $ \longrightarrow $ solution polynomiale
    \item[--] $b(x) = e^{\alpha x} $ $ \longrightarrow $ solution de la forme $Ae^{\alpha x}$
    \item[--] Sinon on utilise la méthode de variation de la constante
\end{itemize}

\begin{equation*}
    ay' + by = b(x)
\end{equation*}



\hrulefill \,
\textbf{Methode de variation de la constante}
\hrulefill

\begin{equation*}
    y_p(x) = \lambda(x)y_h(x)
\end{equation*}

\begin{equation*}
    \lambda'(x) = \frac{b(x)}{y_h(x)}
\end{equation*}

\begin{equation*}
    \lambda(x) = \int \frac{b(x)}{y_h(x)}
\end{equation*}


\hrulefill \,
\textbf{Ordre 2}
\hrulefill 

\end{document}


